Prolog foi inventada em Marseilles (na França), em 1972. Philippe Roussel
escolheu o nome como uma abreviação para ``\textbf{pro}gramação
\textbf{lóg}ica'' \footnote{em francês: ``\textbf{pro}grammation en
\textbf{log}ique''.}, por sugestão de sua esposa no outono de 1972. Pode-se
dizer que Prolog foi a prole de um casamento bem-sucedido entre processamento de
linguagem natural e prova de teorema automatizada. A ideia de usar uma linguagem
natural como o francês para raciocinar e comunicar diretamente com um computador
parecia uma ideia \textit{maluca}, apesar disso, foi a base para o projeto
definido por Alain Colmerauer no verão da década de 1970. \cite{Colmerauer1993-cj}

Notavelmente, Prolog foi peça central do IBM Watson, que em novembro de 2013,
venceu os melhores da categoria no jogo televisivo \textit{Jeopardy!}, sem
conexão com a internet, e precede em 4 anos os modelos de IA baseados em
\textit{transformers} (aprendizado profundo). \cite{watson-jeopardy}

\section{Philippe Roussel}
% por algum motivo, é mais difícil encontrar informações dele, a Wikipédia em
% espanhol tem algumas, mas sem citar fontes.
Philippe Roussel 

\section{Alain Colmerauer}
Alain Colmerauer, foi um Doutor em ciência da computação, professor na
Universidade de Aix-Marselha, e co-criador do Prolog, nascido em 24 de janeiro
de 1941 em Carcassonne. Em 2000, a sua empresa de software se especializava em
análise e gerência de riscos financeiros e hipoteca. E em 2005, foi comprada
pela Experiean, que se tornou parte do negócio de suporte a decisões
empresariais, Experian-Scorex. \cite{Alain-Memorian}

Prêmios recebidos:
\begin{itemize}
    \item Pomme d'Or du Logiciel Francais (1982)
    \item Michel Monpetit Award (1986)
    \item Chevalier de la Legion d'Honneur (1986)
    \item ACP Research Excellence Award (1991)
\end{itemize}


% pesquisar mais afundo acerca do Experian-Scorex